\subsection{Sprint 1}

A Sprint 1 teve como objetivo produzir uma primeira versão navegável e demonstrável do ReadRace. O planejamento apresentado aos clientes previa 37 tarefas, envolvendo a base mobile e backend, componentes reutilizáveis, biblioteca, registro de progresso, funcionalidades sociais, busca, perfil de outro usuário e desafios entre amigos.

Minha atuação combinou desenvolvimento, revisão e gestão. Em 8 de setembro, revisei o PR de migration e seed, participei do alinhamento diário e trabalhei no componente de capa de livro. Em 10 de setembro, concluí atividades relacionadas a BookCover e EmptyState, revisei a tarefa de busca de usuários, livros e comunidades e avancei no desenvolvimento do perfil de usuários.

Nos dias seguintes, participei da resolução de conflitos e da integração de entregas de colegas. Revisei PRs ligados à API de desafios e à biblioteca, levantando a necessidade de avaliar paginação nas listagens. Também trabalhei no componente BottomSheet, no cálculo de XP associado ao registro de leitura, no componente Card, na tela de detalhe do livro e no componente ListItem. Na reta final, desenvolvi e ajustei a API e a tela de visualização do perfil de outro usuário.

A consulta ao GitHub identificou oito PRs de minha autoria, sob o usuário \texttt{Tchuri17}, integrados até 18 de setembro de 2026: BookCover (\#73), EmptyState (\#74), BottomSheet (\#80), cálculo de XP da leitura (\#83), Card (\#84), detalhe do livro (\#86), ListItem (\#88) e perfil de outro usuário (\#91). Os registros estão disponíveis na área de pull requests de \url{https://github.com/ReadRace-AGES/ReadRace}. Esse levantamento complementa os registros pessoais, mas não mede todo o trabalho de gestão, revisão e apoio realizado.

A implementação do registro de leitura exigiu atenção à diferença entre a página atual e a página máxima já alcançada. A regra de XP precisava impedir que a mesma página fosse contabilizada novamente ao retroceder e avançar na leitura. Essa contribuição não corresponde à gamificação completa: a integração com nível, sequência de dias, conquistas e placar de desafios ficou para a etapa seguinte.

Durante as aulas, mantive uma postura solícita, procurando acompanhar o que os colegas estavam fazendo, perguntar sobre o andamento das tarefas e oferecer ajuda. Fora da sala, também entrei em contato para verificar entregas, apoiar a entrada de colegas no GitHub e lembrar a equipe das revisões pendentes. As mensagens de acompanhamento registram pedidos de atualização, avisos sobre o code freeze e a importância de concluir tarefas que liberavam o trabalho de outras pessoas.

Apesar dessas iniciativas, reconheço que nós, AGES IV, deveríamos ter acompanhado o projeto de maneira mais contínua. Faltou cobrar atualizações com antecedência e manter uma visão mais próxima das dificuldades de cada integrante. Na aproximação da entrega, precisei assumir várias tarefas para reduzir as pendências. Minha disponibilidade para ajudar foi positiva, mas a concentração de implementação no final também evidenciou uma falha no acompanhamento e na distribuição do trabalho.

Em 12 de setembro, discutimos o estado do projeto e um plano para as tarefas críticas, buscando contato com seus responsáveis e avaliando a necessidade de assumir entregas sem perspectiva de conclusão no prazo. Em 15 de setembro, trabalhei em ajustes e conflitos relacionados ao detalhamento do livro e reforcei a comunicação sobre a weekly e o encerramento do desenvolvimento. Em 16 de setembro, alinhamos que o fluxo de desafios entre amigos permaneceria pendente, enquanto a equipe concentraria esforços nas demais funcionalidades e na apresentação. No dia seguinte, continuei os ajustes do perfil público e participei da organização da próxima sprint.

A apresentação de resultados de 18 de setembro registrou 32 tarefas entregues, oito telas, 17 endpoints e 16 componentes. Esses números representam o resultado coletivo apresentado aos stakeholders, não entregas individuais minhas. Foram demonstrados biblioteca, detalhe e registro de progresso, feed de comunidades, página do clube com ranking, fórum, busca e perfil de outro leitor. Cinco tarefas do fluxo de desafios entre amigos foram transferidas para a Sprint 2. A interface social dessa etapa era principalmente de consulta; curtidas e outras interações ainda dependiam de evolução.

Na retrospectiva, expliquei que os bloqueios precisam ser comunicados e tratados rapidamente. Uma tarefa parada pode impedir outras entregas e comprometer o projeto inteiro, especialmente nos fluxos com dependências encadeadas. A principal lição foi que a gestão precisa identificar esses sinais durante a sprint, e não apenas quando o prazo está próximo.

Para melhorar esse acompanhamento na Sprint 2, dividimos o time em squads com um ou dois integrantes de AGES IV em cada grupo. O objetivo foi estabelecer responsáveis mais próximos dos desenvolvedores, facilitar a identificação de impedimentos e acompanhar melhor as entregas. A efetividade dessa mudança deverá ser avaliada nas próximas retrospectivas.

Por fim, considero relevante minha participação na apresentação aos stakeholders ao lado dos colegas de gestão. Embora a timidez ainda seja um desafio, preparar o conteúdo e participar da exposição dos resultados tem me permitido praticar uma comunicação mais clara. Nesta sprint, o desenvolvimento pessoal esteve ligado tanto à responsabilidade pelas entregas quanto à capacidade de explicar o que foi realizado e o que permaneceu pendente.
